\section{Multiplicity and typical cardinality}\label{sec:multiplicity}

Fix squarefree $b\ge3$ and retain the direct family $\mathcal E=\mathcal E_X$. Put
\[
M=|\mathcal E|,
\qquad
\mathcal R_X(b)=
\left\{A\subseteq\mathcal E:
\sum_{e\in A}\frac1e=\frac1b\right\}.
\tag{7.1}
\]
Let
\[
q_X=\max(\theta,1-\theta),
\qquad
q_b=\max(\alpha_b,1-\alpha_b).
\tag{7.2}
\]
Then $q_X\to q_b$, while Lemma \ref{lem:family-size} and \eqref{2.19} give
\[
\log L=o(M),
\qquad
|\log\sigma|=o(M).
\tag{7.3}
\]

\subsection{Exponentially many exact representations}

The exact target probability is
\[
P
=
\sum_{A\in\mathcal R_X(b)}
\theta^{|A|}(1-\theta)^{M-|A|}.
\tag{7.4}
\]
Every summand is at most $q_X^M$. Hence
\[
|\mathcal R_X(b)|
\ge Pq_X^{-M}.
\tag{7.5}
\]
By \eqref{6.28},
\[
P\asymp_b\frac1{L\sigma}.
\tag{7.6}
\]
Taking logarithms and using \eqref{7.3},
\[
\log|\mathcal R_X(b)|
\ge
-M\log q_X-o(M)
=
[-\log q_b+o_b(1)]M.
\tag{7.7}
\]
This proves \eqref{1.9}.

\subsection{The typical-cardinality window}

Let
\[
K=\sum_{e\in\mathcal E}\xi_e
\]
be the number of selected denominators under the ambient Bernoulli law, and set
\[
w=
\left\lceil2\sqrt{M\log L}\right\rceil.
\tag{7.8}
\]
Since $\log L=o(M)$,
\[
\sqrt M\ll w=o(M).
\tag{7.9}
\]
Hoeffding's inequality \cite{Hoeffding1963} gives
\[
\Pr(|K-\theta M|>w)
\le
2\exp\left(-\frac{2w^2}{M}\right)
\le
2L^{-8}.
\tag{7.10}
\]
The local limit \eqref{6.28} implies
\[
2L^{-8}=o(P).
\tag{7.11}
\]
Consequently the exact target mass in the window satisfies
\[
\sum_{\substack{A\in\mathcal R_X(b)\\
||A|-\theta M|\le w}}
\theta^{|A|}(1-\theta)^{M-|A|}
=(1-o(1))P.
\tag{7.12}
\]

If $|k-\theta M|\le w$, write $k=\theta M+\delta$. Then
\[
\log\bigl(\theta^k(1-\theta)^{M-k}\bigr)
=
-MH(\theta)
+\delta\log\frac{\theta}{1-\theta},
\tag{7.13}
\]
where
\[
H(u)=-u\log u-(1-u)\log(1-u).
\]
Since $\theta$ remains in a compact subinterval of $(0,1)$,
\[
\theta^k(1-\theta)^{M-k}
\le
\exp\bigl(-MH(\theta)+O_b(w)\bigr)
\tag{7.14}
\]
uniformly in the window. Dividing \eqref{7.12} by this maximum atom and using
\[
\log P=-\log L-\log\sigma+O_b(1)=o(M),
\qquad
w=o(M),
\qquad
H(\theta)\to H(\alpha_b),
\]
gives
\[
\#\left\{A\in\mathcal R_X(b):
\bigl||A|-\theta M\bigr|\le w\right\}
\ge
\exp\bigl([H(\alpha_b)-o_b(1)]M\bigr).
\tag{7.15}
\]
This is \eqref{1.10}.

The exponent is optimal at ambient exponential scale: the standard entropy estimate for binomial coefficients gives
\[
\#\left\{A\subseteq\mathcal E:
\bigl||A|-\theta M\bigr|\le w\right\}
\le
\exp\bigl([H(\alpha_b)+o_b(1)]M\bigr).
\tag{7.16}
\]
Thus imposing the exact reciprocal-sum constraint costs only a subexponential factor within this mesoscopic cardinality window.

\subsection{One exact cardinality}

The interval $|k-\theta M|\le w$ contains at most $2w+1=\exp(o(M))$ integers. By pigeonhole applied to \eqref{7.15}, there is an integer
\[
k_X=\theta M+O(w)
\tag{7.17}
\]
such that
\[
\#\{A\in\mathcal R_X(b):|A|=k_X\}
\ge
\exp\bigl([H(\alpha_b)-o_b(1)]M\bigr).
\tag{7.18}
\]
This proves \eqref{1.11} and completes Theorem \ref{thm:multiplicity}.

The conclusion is deliberately mesoscopic. We do not prove a prescribed-cardinality asymptotic, a conditional central limit theorem for $|A|$ on the natural $\sqrt M$ scale, or a bivariate local limit for reciprocal sum and cardinality.
